\chapter{Backup}
\label{ch:backup}

\section{Backup}\label{backup}

You will find most of the ideas in {[}@ism-book{]}.

\subsection{Terminology}\label{terminology}

\begin{description}
\item[Production data]
is actively used in the course of day-to-day operations.
\item[Backup:]
additional copy of production data that is created and retained for the
sole purpose of recovering lost or corrupted data.
\item[Archive:]
store of data that is no longer actively used that is retained on
low-cost secondary storage.
\end{description}

\subsection{Purpose}\label{purpose}

\begin{description}
\item[Disaster recovery]
where production systems are damaged / destroyed.
\item[Operational recovery]
where data loss / corruption occurs due to application failure or user
error (e.g.~important email deleted).
\item[Archival]
for regulatory compliance, historical record, later analysis etc.
\end{description}

\subsection{Parameters}\label{parameters}

\begin{description}
\item[Recovery Point Objective (RPO)]
the most recent point in time to which we should be able to restore.
Dictates backup frequency.
\item[Recovery Time Objective (RTO)]
is the time taken to restore the system to the RPO point after an
incident occurs.
\end{description}

\subsection{Methods}\label{methods}

\begin{description}
\item[Hot / online:]
production system is running during backup:

\begin{itemize}
\item
  Does not interrupt normal operation
\item
  Issues with open files, particularly database backend stores.
\end{itemize}
\item[Cold / offline:]
production system is shut-down for normal operations.
\end{description}

\subsection{Scope}\label{scope}

Scope refers to what is actually backed up:

\begin{itemize}
\item
  Files themselves
\item
  Files and their layout (e.g.~folders, symlinks)
\item
  Metadata (permissions, last updated)
\item
  Is a \emph{Bare Metal Recovery} required?
\end{itemize}

\section{Backup granularity}\label{backup-granularity}

There are three common granularities of backup, .
\href{https://www.acronis.com/en-us/articles/incremental-differential-backups/}{See
Acronis article} on different granularities for further info.

\subsection{Full backup}\label{full-backup}

Full backup contains the entire dataset:

\begin{itemize}
\item
  Completeness.
\item
  Frequent full backups may be prohibitive in terms of time and/or
  storage space.
\end{itemize}

\subsection{Incremenetal backup}\label{incremenetal-backup}

Incremental backups store only the files changed since the most recent
full or incremental backup was taken. Restoration steps are shown in .

\subsection{Cumulative / Differential
backup}\label{cumulative-differential-backup}

Differential backup includes files changed since the last \textbf{full}
backup. Restoration steps are shown in

\section{Targets}\label{targets}

Backup data is written to targets. Like any other storage media, targets
may be directly attached to a host or may be LAN or SAN attached.

\subsection{Disk}\label{disk}

Disk-based backup has low cost, fast backup and recovery. Should be
using monitored RAID to ensure backup availability.

\subsection{Tape}\label{tape}

Backup traditionally has used tape media.

Tape drives are best operated such that they do not have to start/stop
repeatedly during operation. Interleaving can be used so that multiple
streams of data can be written to a tape at speed, .

\subsection{Tape libraries}\label{tape-libraries}

Tape backup can be provisioned on a tape library in medium-large
installations, .

\subsection{Virtual tape library}\label{virtual-tape-library}

Virtual tape library emulates a tape library but is actually provisioned
usually from SAN LUNs, .

\subsection{Targets comparison}\label{targets-comparison}

Comparison of different targets is given in .
